% GENERATED FILE. Edit canonical lessons, then run npm run generate:books.
% canonical-source-hash: fnv1a64:a0422244e4b3966a
% canonical-lessons: KA-C467-marukatte, KA-C467-pete, KA-C467-malige, KA-C467-marata, KA-C467-konduko

\chapter{At the Market}
\label{ch:ka-at-the-market}
\hlchaptermodality{467}

\begin{chapteropening}
\textbf{By the end of this chapter:} \emph{I can say a marketplace, a bazaar, a store, a sale and to buy for oneself.}

Complete the last lesson of chapter 467: I can say a marketplace, a bazaar, a store, a sale and to buy for oneself.
\end{chapteropening}

\section[\texorpdfstring{\kn{ಮಾರುಕಟ್ಟೆ} (mārukaṭṭe)}{ಮಾರುಕಟ್ಟೆ (mārukaṭṭe)}]{\kn{ಮಾರುಕಟ್ಟೆ} (mārukaṭṭe) — a marketplace; a market}

\label{lesson:KA-C467-marukatte}



\begin{quote}
Before the new one: say the Kannada for a streetlight, then the Kannada for a drain.
\end{quote}

\subsection*{You'll want to know: \kn{ಮಾರುಕಟ್ಟೆ}}
\textbf{\kn{ಮಾರುಕಟ್ಟೆ}} — \emph{mārukaṭṭe} — \textquotedblleft{}a marketplace; a market\textquotedblright{}.

\begin{quote}
\kn{ಭಾನುವಾರ} \kn{ಮಾರುಕಟ್ಟೆಯಲ್ಲಿ} \kn{ತುಂಬಾ} \kn{ಜನ} \kn{ಇರುತ್ತಾರೆ}.
\end{quote}

\begin{quote}
\emph{bhānuvāra mārukaṭṭeyalli tumbā jana iruttāre.}
\end{quote}

\begin{quote}
On Sundays there are a lot of people in the market.
\end{quote}

\kn{ಸಂತೆ} is a weekly fair; \kn{ಮಾರುಕಟ್ಟೆ} is a standing market.

\begin{cousinweb}
\textbf{\kn{ಮಾರುಕಟ್ಟೆ}} (\emph{mārukaṭṭe}) is a native Kannada compound of \emph{māṟu}, \textquotedblleft{}to sell, to barter\textquotedblright{} (DEDR 4834), and \kn{ಕಟ್ಟೆ}, \textquotedblleft{}a raised platform\textquotedblright{}.
\end{cousinweb}

\begin{grammarlens}[title={What you've built}]
One more word for getting around.
\end{grammarlens}

\subsection*{Your turn}
\begin{itemize}
\raggedright
  \item \emph{Say it:} \emph{mārukaṭṭe}
  \item \emph{Say it:} \emph{mārukaṭṭe}, once more
  \item \emph{Say it:} \emph{bhānuvāra mārukaṭṭeyalli tumbā jana iruttāre.} — On Sundays there are a lot of people in the market.
  \item \emph{Recall:} say the Kannada for a streetlight, then the Kannada for a drain, then say \emph{mārukaṭṭe} again
\end{itemize}

\begin{tcolorbox}[breakable,colback=teal!4,colframe=teal!35!black,title={Before you move on}]
What does \kn{ಮಾರುಕಟ್ಟೆ} mean? (\textquotedblleft{}A marketplace; a market\textquotedblright{}.) Say it once more.
\end{tcolorbox}
\section[\texorpdfstring{\kn{ಪೇಟೆ} (pēṭe)}{ಪೇಟೆ (pēṭe)}]{\kn{ಪೇಟೆ} (pēṭe) — a bazaar; a market street}

\label{lesson:KA-C467-pete}



\begin{quote}
Before the new one: say the Kannada for a drain, then the Kannada for a marketplace.
\end{quote}

\subsection*{You'll want to know: \kn{ಪೇಟೆ}}
\textbf{\kn{ಪೇಟೆ}} — \emph{pēṭe} — \textquotedblleft{}a bazaar; a market street\textquotedblright{}.

\begin{quote}
\kn{ಹಳೆಯ} \kn{ಪೇಟೆಯಲ್ಲಿ} \kn{ಬಟ್ಟೆ} \kn{ಅಗ್ಗವಾಗಿ} \kn{ಸಿಗುತ್ತದೆ}.
\end{quote}

\begin{quote}
\emph{haḷeya pēṭeyalli baṭṭe aggavāgi siguttade.}
\end{quote}

\begin{quote}
Cloth is cheap in the old bazaar.
\end{quote}

Also the market quarter of a town: \kn{ಪೇಟೆಗೆ} \kn{ಹೋಗು}, to go into town.

\begin{cousinweb}
\textbf{\kn{ಪೇಟೆ}} (\emph{pēṭe}) is a Kannada word for a market street whose origin is not settled.
\end{cousinweb}

\begin{grammarlens}[title={What you've built}]
One more word for getting around.
\end{grammarlens}

\subsection*{Your turn}
\begin{itemize}
\raggedright
  \item \emph{Say it:} \emph{pēṭe}
  \item \emph{Say it:} \emph{pēṭe}, once more
  \item \emph{Say it:} \emph{haḷeya pēṭeyalli baṭṭe aggavāgi siguttade.} — Cloth is cheap in the old bazaar.
  \item \emph{Recall:} say the Kannada for a drain, then the Kannada for a marketplace, then say \emph{pēṭe} again
\end{itemize}

\begin{tcolorbox}[breakable,colback=teal!4,colframe=teal!35!black,title={Before you move on}]
What does \kn{ಪೇಟೆ} mean? (\textquotedblleft{}A bazaar; a market street\textquotedblright{}.) Say it once more.
\end{tcolorbox}
\section[\texorpdfstring{\kn{ಮಳಿಗೆ} (maḷige)}{ಮಳಿಗೆ (maḷige)}]{\kn{ಮಳಿಗೆ} (maḷige) — a store; a shop unit}

\label{lesson:KA-C467-malige}



\begin{quote}
Before the new one: say the Kannada for a marketplace, then the Kannada for a bazaar.
\end{quote}

\subsection*{You'll want to know: \kn{ಮಳಿಗೆ}}
\textbf{\kn{ಮಳಿಗೆ}} — \emph{maḷige} — \textquotedblleft{}a store; a shop unit\textquotedblright{}.

\begin{quote}
\kn{ಈ} \kn{ಬೀದಿಯಲ್ಲಿ} \kn{ಬಟ್ಟೆಯ} \kn{ಮಳಿಗೆಗಳು} \kn{ಇವೆ}.
\end{quote}

\begin{quote}
\emph{ī bīdiyalli baṭṭeya maḷigegaḷu ive.}
\end{quote}

\begin{quote}
There are clothes stores on this street.
\end{quote}

Plural: \kn{ಮಳಿಗೆಗಳು}, \emph{maḷigegaḷu}.

A large shop or a showroom; \kn{ಅಂಗಡಿ} is any shop.

\begin{cousinweb}
\textbf{\kn{ಮಳಿಗೆ}} (\emph{maḷige}) is native Dravidian: DEDR 4757 lists Kannada \emph{maḷige}, \textquotedblleft{}a warehouse, a shop\textquotedblright{}.
\end{cousinweb}

\begin{grammarlens}[title={What you've built}]
One more word for getting around.
\end{grammarlens}

\subsection*{Your turn}
\begin{itemize}
\raggedright
  \item \emph{Say it:} \emph{maḷige}
  \item \emph{Say it:} \emph{maḷige}, once more
  \item \emph{Say it:} \emph{ī bīdiyalli baṭṭeya maḷigegaḷu ive.} — There are clothes stores on this street.
  \item \emph{Recall:} say the Kannada for a marketplace, then the Kannada for a bazaar, then say \emph{maḷige} again
\end{itemize}

\begin{tcolorbox}[breakable,colback=teal!4,colframe=teal!35!black,title={Before you move on}]
What does \kn{ಮಳಿಗೆ} mean? (\textquotedblleft{}A store; a shop unit\textquotedblright{}.) Say it once more.
\end{tcolorbox}
\section[\texorpdfstring{\kn{ಮಾರಾಟ} (mārāṭa)}{ಮಾರಾಟ (mārāṭa)}]{\kn{ಮಾರಾಟ} (mārāṭa) — a sale; selling}

\label{lesson:KA-C467-marata}



\begin{quote}
Before the new one: say the Kannada for a bazaar, then the Kannada for a store.
\end{quote}

\subsection*{You'll want to know: \kn{ಮಾರಾಟ}}
\textbf{\kn{ಮಾರಾಟ}} — \emph{mārāṭa} — \textquotedblleft{}a sale; selling\textquotedblright{}.

\begin{quote}
\kn{ಹಬ್ಬದ} \kn{ಮುಂಚೆ} \kn{ಸೀರೆಗಳ} \kn{ಮಾರಾಟ} \kn{ಹೆಚ್ಚಾಗುತ್ತದೆ}.
\end{quote}

\begin{quote}
\emph{habbada muñce sīregaḷa mārāṭa heccāguttade.}
\end{quote}

\begin{quote}
Before a festival the sale of saris goes up.
\end{quote}

\kn{ಮಾರಾಟಕ್ಕೆ} \kn{ಇದೆ}, it is for sale; \kn{ಮಾರಾಟಗಾರ}, a seller.

\begin{cousinweb}
\textbf{\kn{ಮಾರಾಟ}} (\emph{mārāṭa}) is native Dravidian, a noun built on \emph{māṟu}, \textquotedblleft{}to sell, to barter\textquotedblright{} (DEDR 4834).
\end{cousinweb}

\begin{grammarlens}[title={What you've built}]
One more word for getting around.
\end{grammarlens}

\subsection*{Your turn}
\begin{itemize}
\raggedright
  \item \emph{Say it:} \emph{mārāṭa}
  \item \emph{Say it:} \emph{mārāṭa}, once more
  \item \emph{Say it:} \emph{habbada muñce sīregaḷa mārāṭa heccāguttade.} — Before a festival the sale of saris goes up.
  \item \emph{Recall:} say the Kannada for a bazaar, then the Kannada for a store, then say \emph{mārāṭa} again
\end{itemize}

\begin{tcolorbox}[breakable,colback=teal!4,colframe=teal!35!black,title={Before you move on}]
What does \kn{ಮಾರಾಟ} mean? (\textquotedblleft{}A sale; selling\textquotedblright{}.) Say it once more.
\end{tcolorbox}
\section[\texorpdfstring{\kn{ಕೊಂಡುಕೊ} (koṇḍuko)}{ಕೊಂಡುಕೊ (koṇḍuko)}]{\kn{ಕೊಂಡುಕೊ} (koṇḍuko) — to buy for oneself; to purchase}

\label{lesson:KA-C467-konduko}



\begin{quote}
Before the new one: say the Kannada for a store, then the Kannada for a sale.
\end{quote}

\subsection*{You'll want to know: \kn{ಕೊಂಡುಕೊ}}
\textbf{\kn{ಕೊಂಡುಕೊ}} — \emph{koṇḍuko} — \textquotedblleft{}to buy for oneself; to purchase\textquotedblright{}.

\begin{quote}
\kn{ನಾನು} \kn{ಮಾರುಕಟ್ಟೆಯಲ್ಲಿ} \kn{ಹೊಸ} \kn{ಚಪ್ಪಲಿ} \kn{ಕೊಂಡುಕೊಂಡೆ}.
\end{quote}

\begin{quote}
\emph{nānu mārukaṭṭeyalli hosa cappali koṇḍukoṇḍe.}
\end{quote}

\begin{quote}
I bought new sandals at the market.
\end{quote}

Past (I …): \kn{ಕೊಂಡುಕೊಂಡೆ}, \emph{koṇḍukoṇḍe}.

\kn{ಕೊಂಡು} is the past stem of \kn{ಕೊಳ್ಳು}, to take or buy; \kn{ಕೊ} adds for oneself.

\begin{cousinweb}
\textbf{\kn{ಕೊಂಡುಕೊ}} (\emph{koṇḍuko}) is native Dravidian: \emph{koṇḍu}, the past stem of \emph{koḷ}, \emph{koḷḷu}, \textquotedblleft{}to take, to obtain, to buy\textquotedblright{} (DEDR 2151), with the reflexive -\emph{ko}.
\end{cousinweb}

\begin{grammarlens}[title={What you've built}]
One more word for getting around.
\end{grammarlens}

\subsection*{Your turn}
\begin{itemize}
\raggedright
  \item \emph{Say it:} \emph{koṇḍuko}
  \item \emph{Say it:} \emph{koṇḍuko}, once more
  \item \emph{Say it:} \emph{nānu mārukaṭṭeyalli hosa cappali koṇḍukoṇḍe.} — I bought new sandals at the market.
  \item \emph{Recall:} say the Kannada for a store, then the Kannada for a sale, then say \emph{koṇḍuko} again
\end{itemize}

\begin{tcolorbox}[breakable,colback=teal!4,colframe=teal!35!black,title={Before you move on}]
What does \kn{ಕೊಂಡುಕೊ} mean? (\textquotedblleft{}To buy for oneself; to purchase\textquotedblright{}.) Say it once more.
\end{tcolorbox}
